\begin{lemma}
\label{lemma-quasi-compact-h1-zero-covering}
\begin{reference}
\cite{Serre-criterion}, \cite[II, Theorem 5.2.1 (d') and IV (1.7.17)]{EGA}
\end{reference}
\begin{slogan}
Serre's criterion for affineness.
\end{slogan}
Let $X$ be a scheme.
Assume that
\begin{enumerate}
\item $X$ is quasi-compact,
\item for every quasi-coherent sheaf of ideals
$\mathcal{I} \subset \mathcal{O}_X$ we have $H^1(X, \mathcal{I}) = 0$.
\end{enumerate}
Then $X$ is affine.
\end{lemma}

\begin{proof}
Let $x \in X$ be a closed point. Let $U \subset X$ be an affine open
neighbourhood of $x$. Write $U = \Spec(A)$ and let
$\mathfrak m \subset A$ be the maximal ideal corresponding to $x$.
Set $Z = X \setminus U$ and $Z' = Z \cup \{x\}$.
By Schemes, Lemma \ref{schemes-lemma-reduced-closed-subscheme} there
are quasi-coherent sheaves of ideals
$\mathcal{I}$, resp.\ $\mathcal{I}'$ cutting out
the reduced closed subschemes $Z$, resp.\ $Z'$.
Consider the short exact sequence
$$
0 \to \mathcal{I}' \to \mathcal{I} \to \mathcal{I}/\mathcal{I}' \to 0.
$$
Since $x$ is a closed point of $X$ and $x \not \in Z$ we see that
$\mathcal{I}/\mathcal{I}'$ is supported at $x$. In fact, the restriction
of $\mathcal{I}/\mathcal{I'}$ to $U$ corresponds to the $A$-module
$A/\mathfrak m$. Hence we see that $\Gamma(X, \mathcal{I}/\mathcal{I'})
= A/\mathfrak m$. Since by assumption $H^1(X, \mathcal{I}') = 0$
we see there exists a global section $f \in \Gamma(X, \mathcal{I})$
which maps to the element $1 \in A/\mathfrak m$ as a section of
$\mathcal{I}/\mathcal{I'}$. Clearly we have
$x \in X_f \subset U$. This implies that $X_f = D(f_A)$ where
$f_A$ is the image of $f$ in $A = \Gamma(U, \mathcal{O}_X)$.
In particular $X_f$ is affine.

\medskip\noindent
Consider the union $W = \bigcup X_f$ over all $f \in \Gamma(X, \mathcal{O}_X)$
such that $X_f$ is affine. Obviously $W$ is open in $X$.
By the arguments above every closed point of
$X$ is contained in $W$. The closed subset $X \setminus W$ of $X$
is also quasi-compact
(see Topology, Lemma \ref{topology-lemma-closed-in-quasi-compact}).
Hence it has a closed point if it is nonempty (see
Topology, Lemma \ref{topology-lemma-quasi-compact-closed-point}).
This would contradict the fact that all closed points are in
$W$. Hence we conclude $X = W$.

\medskip\noindent
Choose finitely many $f_1, \ldots, f_n \in \Gamma(X, \mathcal{O}_X)$
such that $X = X_{f_1} \cup \ldots \cup X_{f_n}$ and such that each
$X_{f_i}$ is affine. This is possible as we've seen above.
By Properties, Lemma \ref{properties-lemma-characterize-affine}
to finish the proof it suffices
to show that $f_1, \ldots, f_n$ generate the unit ideal in
$\Gamma(X, \mathcal{O}_X)$. Consider the short exact sequence
$$
\xymatrix{
0 \ar[r] &
\mathcal{F} \ar[r] &
\mathcal{O}_X^{\oplus n} \ar[rr]^{f_1, \ldots, f_n} & &
\mathcal{O}_X \ar[r] &
0
}
$$
The arrow defined by $f_1, \ldots, f_n$ is surjective since the
opens $X_{f_i}$ cover $X$. We let $\mathcal{F}$ be the kernel
of this surjective map.
Observe that $\mathcal{F}$ has a filtration
$$
0 = \mathcal{F}_0 \subset \mathcal{F}_1 \subset
\ldots \subset \mathcal{F}_n = \mathcal{F}
$$
so that each subquotient $\mathcal{F}_i/\mathcal{F}_{i - 1}$ is
isomorphic to a quasi-coherent sheaf of ideals.
Namely we can take $\mathcal{F}_i$ to be the intersection of
$\mathcal{F}$ with the first $i$ direct summands of
$\mathcal{O}_X^{\oplus n}$.
The assumption
of the lemma implies that $H^1(X, \mathcal{F}_i/\mathcal{F}_{i - 1}) = 0$
for all $i$. This implies that
$H^1(X, \mathcal{F}_2) = 0$ because it is sandwiched between
$H^1(X, \mathcal{F}_1)$ and $H^1(X, \mathcal{F}_2/\mathcal{F}_1)$.
Continuing like this we deduce that $H^1(X, \mathcal{F}) = 0$.
Therefore we conclude that the map
$$
\xymatrix{
\bigoplus\nolimits_{i = 1, \ldots, n} \Gamma(X, \mathcal{O}_X)
\ar[rr]^{f_1, \ldots, f_n} & &
\Gamma(X, \mathcal{O}_X)
}
$$
is surjective as desired.
\end{proof}
